
\documentclass[10pt,letterpaper]{article}
\usepackage[T1]{fontenc}
\usepackage[utf8]{inputenc}
\usepackage{newtxtext}
\usepackage[margin=0.72in,headheight=14pt]{geometry}
\usepackage{xcolor}
\usepackage{hyperref}
\usepackage{bookmark}
\usepackage{enumitem}
\usepackage{booktabs}
\usepackage{longtable}
\usepackage{array}
\usepackage{ragged2e}
\usepackage{titlesec}
\usepackage{fancyhdr}
\usepackage{microtype}
\usepackage{listings}
\usepackage{needspace}
\usepackage{tocloft}
\usepackage{upquote}
\definecolor{HouseBlue}{HTML}{1F4E79}
\definecolor{HouseGray}{HTML}{4F5B62}
\definecolor{HouseLight}{HTML}{EAF2F8}
\definecolor{CodeBack}{HTML}{F7F8FA}
\definecolor{CodeFrame}{HTML}{C9D1D9}
\definecolor{CodeKeyword}{HTML}{1F4E79}
\definecolor{CodeComment}{HTML}{5F6F7A}
\definecolor{CodeString}{HTML}{8A3B12}
\hypersetup{colorlinks=true,linkcolor=HouseBlue,urlcolor=HouseBlue,citecolor=HouseBlue,pdfauthor={ChatGPT},pdftitle={Game Programming Patterns: Algorithms as C-like Pseudocode}}
\pagestyle{fancy}
\fancyhf{}
\lhead{\small Game Programming Patterns --- Algorithms as C-like Pseudocode}
\rhead{\small\thepage}
\renewcommand{\headrulewidth}{0.4pt}
\setlength{\parindent}{0pt}
\setlength{\parskip}{5pt plus 1pt}
\setlist[itemize]{leftmargin=1.25em,itemsep=2pt,topsep=2pt}
\setlist[enumerate]{leftmargin=1.5em,itemsep=2pt,topsep=2pt}
\titleformat{\section}{\Large\bfseries\color{HouseBlue}}{}{0pt}{}
\titleformat{\subsection}{\large\bfseries\color{HouseGray}}{}{0pt}{}
\titleformat{\subsubsection}{\normalsize\bfseries\color{HouseGray}}{}{0pt}{}
\titlespacing*{\section}{0pt}{2.2ex plus .4ex minus .2ex}{1.0ex}
\titlespacing*{\subsection}{0pt}{1.8ex plus .3ex minus .2ex}{0.7ex}
\newcolumntype{L}[1]{>{\RaggedRight\arraybackslash}p{#1}}
\lstdefinelanguage{PseudoC}{
  morekeywords={typedef,enum,struct,union,return,if,else,for,while,do,case,switch,break,continue,bool,true,false,TRUE,FALSE,NULL,void,int,long,char,float,double,string,const,static,inline,foreach,in,new,delete,insert,remove,push,pop,emit,try,catch,throw,private,public,Command,Actor,Entity,GameObject,Component,Event,EventType,State,StateId,StateStack,Service,Audio,Graphics,Physics,Input,Transform,Vec2,Vec3,Color,Texture,Mesh,Queue,Stack,List,Set,Map,Array,Grid,Cell,Unit,Particle,Breed,Monster,Subject,Observer,Prototype,Pool,Node,Bytecode,Instruction,VM,Frame,Framebuffer,World,Scene,GameContext},
  sensitive=true,
  morecomment=[l]{//},
  morecomment=[s]{/*}{*/},
  morestring=[b]",
}
\lstset{
  language=PseudoC,
  basicstyle=\ttfamily\scriptsize,
  keywordstyle=\bfseries\color{CodeKeyword},
  commentstyle=\itshape\color{CodeComment},
  stringstyle=\color{CodeString},
  backgroundcolor=\color{CodeBack},
  frame=single,
  rulecolor=\color{CodeFrame},
  framerule=0.35pt,
  breaklines=true,
  breakatwhitespace=false,
  columns=fullflexible,
  keepspaces=true,
  showstringspaces=false,
  tabsize=4,
  xleftmargin=0.6em,
  xrightmargin=0.3em,
  aboveskip=0.7em,
  belowskip=0.7em,
}
\newcommand{\handouttitle}{Game Programming Patterns: Algorithms as C-like Pseudocode}
\sloppy
\begin{document}
\pagenumbering{gobble}
\begin{titlepage}
\centering
\vspace*{0.82in}
{\Huge\bfseries\color{HouseBlue}\handouttitle\par}
\vspace{0.20in}
{\Large Systematic Pattern Algorithm Extraction and Reconstruction\par}
\vspace{0.32in}
{\large Source analyzed: Bob Nystrom, Game Programming Patterns, 2014\par}
\vfill
\begin{minipage}{0.88\textwidth}
\small
This handout rewrites the algorithmic behavior expressed in the uploaded game-programming-patterns text as independent C-like pseudocode. It is organized as an implementation-oriented reference for command dispatch, flyweight sharing, observer notification, prototype-based construction, state machines, double buffering, game loops, bytecode interpreters, component systems, event queues, service location, cache-aware data layout, dirty flags, object pools, and spatial partitions. The pseudocode is normalized for study, code review, engine design, and implementation planning; it is not a source-code clone.
\end{minipage}
\vfill
{\large Generated \today\par}
\end{titlepage}
\clearpage
\pagenumbering{roman}
\tableofcontents
\clearpage
\pagenumbering{arabic}

\section*{Extraction Coverage}
\addcontentsline{toc}{section}{Extraction Coverage}
\begin{longtable}{L{0.27\textwidth}L{0.67\textwidth}}
\toprule
\textbf{Area} & \textbf{Algorithm families represented}\\
\midrule
Behavior reification & Command binding, actor-parameterized commands, AI command streams, undoable commands, redo history.\\
Data sharing and construction & Flyweight intrinsic/extrinsic state, terrain flyweights, lazy factories, prototype cloning, spawn callbacks, prototype-data delegation.\\
Runtime communication & Observer lists, intrusive observers, safe destruction, event queues, audio request coalescing, component messages.\\
State and sequencing & Enum FSMs, state object dispatch, local state timers, concurrent FSMs, hierarchical FSMs, pushdown automata, double buffering, game loops, update methods.\\
Scripted behavior & Stack VM bytecode execution, bytecode disassembly, subclass sandbox operations, type-object breed inheritance.\\
Engine structure & Component delegation, service locator registration and lookup, null services, logged service decorators.\\
Performance patterns & Packed arrays, active-prefix partitioning, hot/cold splitting, dirty-flag invalidation, object pools, intrusive grid cells, nearby-interaction queries.\\
\bottomrule
\end{longtable}

\textbf{Source analyzed:} Bob Nystrom, \emph{Game Programming Patterns}, 2014. \textbf{Output purpose:} implementation-oriented C-like pseudocode reference for game architecture and runtime patterns. \textbf{Method:} the book's C++ examples, diagrams, and pattern implementation descriptions have been reduced to independent C-like pseudocode. Names and control flow are normalized for readability and implementation planning.

\medskip\hrule\medskip


\textbf{Source analyzed:} \emph{Game Programming Patterns} by Bob Nystrom, uploaded PDF.
\textbf{Scope:} This handout systematically extracts the algorithmic behavior expressed in the book's source-code examples and rewrites it as implementation-neutral C-like pseudocode. It is not a verbatim code copy. It normalizes names, removes language-specific syntax, and emphasizes control flow, state transitions, data ownership, and runtime behavior.

\Needspace{8\baselineskip}
\subsection{Algorithm Catalog}

\begin{longtable}{L{0.20\textwidth}L{0.34\textwidth}L{0.36\textwidth}}
\toprule
\textbf{Section} & \textbf{Algorithm Family} & \textbf{Core Idea} \\
\midrule\endfirsthead
\toprule
\textbf{Section} & \textbf{Algorithm Family} & \textbf{Core Idea} \\
\midrule\endhead
Command & Input command binding, actor commands, undo/redo & Reify actions as command objects or records. \\
Flyweight & Intrinsic/extrinsic state sharing & Store shared immutable data once and reference it from many instances. \\
Observer & Synchronous notification, linked observer lists & Subjects notify decoupled observers through a stable callback interface. \\
Prototype & Clone-based spawning, spawn callbacks, data delegation & Create new objects from exemplar objects or prototype data. \\
Singleton & Controlled single instance, access alternatives & Enforce one instance or provide service access without over-coupling. \\
State & FSMs, state objects, hierarchical and pushdown state machines & Replace flag tangles with explicit state transition logic. \\
Double Buffer & Framebuffer swap, buffered actor state & Write to next state while readers see current state, then swap atomically. \\
Game Loop & Variable timestep, fixed timestep, accumulator loop & Advance simulation and rendering under real-time constraints. \\
Update Method & Per-frame object simulation & Let each object advance one frame of behavior at a time. \\
Bytecode & Stack VM interpreter & Execute compact instructions using an instruction pointer and operand stack. \\
Subclass Sandbox & Restricted high-level primitive API & Let subclasses define behavior only through safe base-class operations. \\
Type Object & Runtime type records, breed inheritance & Move type-specific data out of classes and into type objects. \\
Component & Input/physics/graphics composition & Compose entities from independently updateable subsystems. \\
Event Queue & Queued communication, audio coalescing & Decouple event production from event processing in time. \\
Service Locator & Global service indirection, null and logged services & Provide globally reachable services without binding users to concrete classes. \\
Data Locality & Packed arrays, hot/cold splitting, active partitioning & Shape memory layout around cache-friendly iteration. \\
Dirty Flag & Lazy derived-data recomputation & Mark derived state invalid and recompute only on demand. \\
Object Pool & Fixed allocation pool, free list & Reuse preallocated objects to avoid allocation churn and fragmentation. \\
Spatial Partition & Uniform grid, linked cells, neighbor queries & Reduce pairwise spatial tests by placing objects into cells. \\
\bottomrule
\end{longtable}

\medskip\hrule\medskip

\Needspace{8\baselineskip}
\subsection{Shared Pseudocode Conventions}

\begin{lstlisting}
#define TRUE 1
#define FALSE 0
#define NULL 0

typedef int Bool;
typedef int Button;
typedef int EventType;
typedef int StateId;
typedef int EntityId;

typedef struct Vec2 { float x, y; } Vec2;
typedef struct Vec3 { float x, y, z; } Vec3;

typedef struct Color { float r, g, b, a; } Color;
\end{lstlisting}

The pseudocode uses explicit structs and function pointers where the original C++ examples rely on classes, inheritance, or virtual dispatch.

\medskip\hrule\medskip

\clearpage
\section{1. Command Pattern Algorithms}

\Needspace{8\baselineskip}
\subsection{1.1 Input-to-Command Dispatch}

\textbf{Purpose:} Replace hardwired input conditionals with configurable command bindings.

\begin{lstlisting}
typedef struct Command Command;
struct Command {
    void (*execute)(Command *self);
};

typedef struct InputBindings {
    Command *button_x;
    Command *button_y;
    Command *button_a;
    Command *button_b;
} InputBindings;

void handle_input(InputBindings *bindings) {
    if (is_pressed(BUTTON_X)) {
        bindings->button_x->execute(bindings->button_x);
    } else if (is_pressed(BUTTON_Y)) {
        bindings->button_y->execute(bindings->button_y);
    } else if (is_pressed(BUTTON_A)) {
        bindings->button_a->execute(bindings->button_a);
    } else if (is_pressed(BUTTON_B)) {
        bindings->button_b->execute(bindings->button_b);
    }
}
\end{lstlisting}

\Needspace{8\baselineskip}
\subsection{1.2 Null Command}

\textbf{Purpose:} Avoid null checks by binding inactive inputs to a do-nothing command.

\begin{lstlisting}
void null_command_execute(Command *self) {
    /* intentionally empty */
}

Command make_null_command(void) {
    Command command;
    command.execute = null_command_execute;
    return command;
}
\end{lstlisting}

\Needspace{8\baselineskip}
\subsection{1.3 Actor-Parameterized Command}

\textbf{Purpose:} Decouple command selection from the actor that performs the command.

\begin{lstlisting}
typedef struct Actor Actor;

typedef struct ActorCommand ActorCommand;
struct ActorCommand {
    void (*execute)(ActorCommand *self, Actor *actor);
};

void jump_command_execute(ActorCommand *self, Actor *actor) {
    actor_jump(actor);
}

ActorCommand *poll_actor_command(InputBindings *bindings) {
    if (is_pressed(BUTTON_X)) return (ActorCommand *)bindings->button_x;
    if (is_pressed(BUTTON_Y)) return (ActorCommand *)bindings->button_y;
    if (is_pressed(BUTTON_A)) return (ActorCommand *)bindings->button_a;
    if (is_pressed(BUTTON_B)) return (ActorCommand *)bindings->button_b;
    return NULL;
}

void dispatch_player_command(InputBindings *bindings, Actor *player) {
    ActorCommand *command = poll_actor_command(bindings);
    if (command != NULL) {
        command->execute(command, player);
    }
}
\end{lstlisting}

\Needspace{8\baselineskip}
\subsection{1.4 AI Command Stream}

\textbf{Purpose:} Let AI and input systems both produce command objects consumed by actors.

\begin{lstlisting}
typedef struct CommandQueue {
    ActorCommand *items[1024];
    int head;
    int tail;
} CommandQueue;

void enqueue_command(CommandQueue *q, ActorCommand *command) {
    q->items[q->tail] = command;
    q->tail = (q->tail + 1) % 1024;
}

ActorCommand *dequeue_command(CommandQueue *q) {
    if (q->head == q->tail) return NULL;
    ActorCommand *command = q->items[q->head];
    q->head = (q->head + 1) % 1024;
    return command;
}

void process_command_stream(CommandQueue *q, Actor *actor) {
    ActorCommand *command;
    while ((command = dequeue_command(q)) != NULL) {
        command->execute(command, actor);
    }
}
\end{lstlisting}

\Needspace{8\baselineskip}
\subsection{1.5 Undoable Move Command}

\textbf{Purpose:} Store enough pre-state to reverse a command after execution.

\begin{lstlisting}
typedef struct Unit Unit;

typedef struct MoveUnitCommand {
    Unit *unit;
    int x_before;
    int y_before;
    int x_after;
    int y_after;
} MoveUnitCommand;

void execute_move_unit(MoveUnitCommand *command) {
    command->x_before = unit_x(command->unit);
    command->y_before = unit_y(command->unit);
    unit_move_to(command->unit, command->x_after, command->y_after);
}

void undo_move_unit(MoveUnitCommand *command) {
    unit_move_to(command->unit, command->x_before, command->y_before);
}
\end{lstlisting}

\Needspace{8\baselineskip}
\subsection{1.6 Undo/Redo History}

\textbf{Purpose:} Maintain a command list and a cursor into the last applied command.

\begin{lstlisting}
#define MAX_HISTORY 4096

typedef struct UndoCommand UndoCommand;
struct UndoCommand {
    void (*execute)(UndoCommand *self);
    void (*undo)(UndoCommand *self);
};

typedef struct CommandHistory {
    UndoCommand *commands[MAX_HISTORY];
    int count;
    int cursor; /* index of next command to redo */
} CommandHistory;

void history_execute(CommandHistory *history, UndoCommand *command) {
    while (history->count > history->cursor) {
        destroy_command(history->commands[history->count - 1]);
        history->count--;
    }

    command->execute(command);
    history->commands[history->count++] = command;
    history->cursor = history->count;
}

void history_undo(CommandHistory *history) {
    if (history->cursor == 0) return;
    history->cursor--;
    history->commands[history->cursor]->undo(history->commands[history->cursor]);
}

void history_redo(CommandHistory *history) {
    if (history->cursor >= history->count) return;
    history->commands[history->cursor]->execute(history->commands[history->cursor]);
    history->cursor++;
}
\end{lstlisting}

\medskip\hrule\medskip

\clearpage
\section{2. Flyweight Pattern Algorithms}

\Needspace{8\baselineskip}
\subsection{2.1 Split Intrinsic and Extrinsic State}

\textbf{Purpose:} Store shared immutable state once and per-instance state separately.

\begin{lstlisting}
typedef struct TreeModel {
    Mesh *mesh;
    Texture *bark;
    Texture *leaves;
} TreeModel;

typedef struct TreeInstance {
    TreeModel *model;      /* intrinsic shared state */
    Vec3 position;         /* extrinsic instance state */
    float height;
    float thickness;
    Color bark_tint;
    Color leaf_tint;
} TreeInstance;

void render_forest(TreeInstance *trees, int count) {
    for (int i = 0; i < count; i++) {
        gpu_draw_instance(
            trees[i].model->mesh,
            trees[i].model->bark,
            trees[i].model->leaves,
            trees[i].position,
            trees[i].height,
            trees[i].thickness,
            trees[i].bark_tint,
            trees[i].leaf_tint);
    }
}
\end{lstlisting}

\Needspace{8\baselineskip}
\subsection{2.2 Terrain Flyweight Grid}

\textbf{Purpose:} Represent each tile by a pointer to immutable terrain data instead of duplicating terrain attributes.

\begin{lstlisting}
#define WORLD_W 256
#define WORLD_H 256

typedef struct Terrain {
    int movement_cost;
    Bool is_water;
    Texture *texture;
} Terrain;

typedef struct World {
    Terrain grass;
    Terrain hill;
    Terrain river;
    Terrain *tiles[WORLD_W][WORLD_H];
} World;

void world_init(World *world) {
    world->grass = (Terrain){ .movement_cost = 1, .is_water = FALSE, .texture = texture_grass() };
    world->hill  = (Terrain){ .movement_cost = 3, .is_water = FALSE, .texture = texture_hill() };
    world->river = (Terrain){ .movement_cost = 2, .is_water = TRUE,  .texture = texture_river() };
}

void generate_terrain(World *world) {
    for (int x = 0; x < WORLD_W; x++) {
        for (int y = 0; y < WORLD_H; y++) {
            if (random_int(10) == 0) world->tiles[x][y] = &world->hill;
            else world->tiles[x][y] = &world->grass;
        }
    }

    int river_x = random_int(WORLD_W);
    for (int y = 0; y < WORLD_H; y++) {
        world->tiles[river_x][y] = &world->river;
    }
}

Terrain *world_tile(World *world, int x, int y) {
    return world->tiles[x][y];
}
\end{lstlisting}

\Needspace{8\baselineskip}
\subsection{2.3 Lazy Flyweight Factory}

\textbf{Purpose:} Reuse an existing flyweight if one with matching intrinsic properties already exists.

\begin{lstlisting}
#define MAX_FLYWEIGHTS 1024

typedef struct FlyweightFactory {
    Terrain items[MAX_FLYWEIGHTS];
    int count;
} FlyweightFactory;

Terrain *get_terrain_flyweight(FlyweightFactory *factory,
                               int movement_cost,
                               Bool is_water,
                               Texture *texture) {
    for (int i = 0; i < factory->count; i++) {
        Terrain *t = &factory->items[i];
        if (t->movement_cost == movement_cost &&
            t->is_water == is_water &&
            t->texture == texture) {
            return t;
        }
    }

    Terrain *created = &factory->items[factory->count++];
    *created = (Terrain){ movement_cost, is_water, texture };
    return created;
}
\end{lstlisting}

\medskip\hrule\medskip

\clearpage
\section{3. Observer Pattern Algorithms}

\Needspace{8\baselineskip}
\subsection{3.1 Array-backed Subject Notification}

\textbf{Purpose:} Let a subject synchronously notify many observers without knowing their concrete types.

\begin{lstlisting}
#define MAX_OBSERVERS 64

typedef struct Entity Entity;

typedef struct Observer Observer;
struct Observer {
    void (*on_notify)(Observer *self, Entity *entity, EventType event);
};

typedef struct Subject {
    Observer *observers[MAX_OBSERVERS];
    int count;
} Subject;

void subject_add(Subject *subject, Observer *observer) {
    subject->observers[subject->count++] = observer;
}

void subject_remove(Subject *subject, Observer *observer) {
    for (int i = 0; i < subject->count; i++) {
        if (subject->observers[i] == observer) {
            subject->observers[i] = subject->observers[subject->count - 1];
            subject->count--;
            return;
        }
    }
}

void subject_notify(Subject *subject, Entity *entity, EventType event) {
    for (int i = 0; i < subject->count; i++) {
        subject->observers[i]->on_notify(subject->observers[i], entity, event);
    }
}
\end{lstlisting}

\Needspace{8\baselineskip}
\subsection{3.2 Achievement Observer}

\textbf{Purpose:} Derive higher-level game achievements from domain events without coupling the domain system to achievements.

\begin{lstlisting}
typedef struct AchievementObserver {
    Observer base;
    Bool hero_was_on_bridge;
} AchievementObserver;

void achievements_on_notify(Observer *base, Entity *entity, EventType event) {
    AchievementObserver *self = (AchievementObserver *)base;

    switch (event) {
        case EVENT_ENTITY_STARTED_FALLING:
            if (entity_is_hero(entity) && self->hero_was_on_bridge) {
                unlock_achievement(ACHIEVEMENT_FELL_OFF_BRIDGE);
            }
            break;

        case EVENT_HERO_ENTERED_BRIDGE:
            self->hero_was_on_bridge = TRUE;
            break;

        case EVENT_HERO_LEFT_BRIDGE:
            self->hero_was_on_bridge = FALSE;
            break;
    }
}
\end{lstlisting}

\Needspace{8\baselineskip}
\subsection{3.3 Intrusive Linked Observer List}

\textbf{Purpose:} Avoid dynamic allocation for the observer list by storing next pointers inside observer records.

\begin{lstlisting}
typedef struct LinkedObserver LinkedObserver;
struct LinkedObserver {
    void (*on_notify)(LinkedObserver *self, Entity *entity, EventType event);
    LinkedObserver *next;
};

typedef struct LinkedSubject {
    LinkedObserver *head;
} LinkedSubject;

void linked_subject_add(LinkedSubject *subject, LinkedObserver *observer) {
    observer->next = subject->head;
    subject->head = observer;
}

void linked_subject_remove(LinkedSubject *subject, LinkedObserver *observer) {
    if (subject->head == observer) {
        subject->head = observer->next;
        observer->next = NULL;
        return;
    }

    LinkedObserver *current = subject->head;
    while (current != NULL) {
        if (current->next == observer) {
            current->next = observer->next;
            observer->next = NULL;
            return;
        }
        current = current->next;
    }
}

void linked_subject_notify(LinkedSubject *subject, Entity *entity, EventType event) {
    LinkedObserver *observer = subject->head;
    while (observer != NULL) {
        observer->on_notify(observer, entity, event);
        observer = observer->next;
    }
}
\end{lstlisting}

\Needspace{8\baselineskip}
\subsection{3.4 Safe Destruction Protocol}

\textbf{Purpose:} Prevent dangling subject-to-observer pointers.

\begin{lstlisting}
void observer_destroy(Observer *observer, Subject **subjects, int subject_count) {
    for (int i = 0; i < subject_count; i++) {
        subject_remove(subjects[i], observer);
    }
    free(observer);
}

void subject_destroy(Subject *subject) {
    subject_notify(subject, NULL, EVENT_SUBJECT_DESTROYING);
    subject->count = 0;
    free(subject);
}
\end{lstlisting}

\medskip\hrule\medskip

\clearpage
\section{4. Prototype Pattern Algorithms}

\Needspace{8\baselineskip}
\subsection{4.1 Clone-based Spawner}

\textbf{Purpose:} Create new instances by cloning a prototype object.

\begin{lstlisting}
typedef struct Monster Monster;
struct Monster {
    Monster *(*clone)(Monster *self);
    int health;
    int speed;
};

typedef struct Spawner {
    Monster *prototype;
} Spawner;

Monster *spawner_spawn(Spawner *spawner) {
    return spawner->prototype->clone(spawner->prototype);
}
\end{lstlisting}

\Needspace{8\baselineskip}
\subsection{4.2 Concrete Clone Routine}

\textbf{Purpose:} Copy class identity and selected state into a new object.

\begin{lstlisting}
Monster *clone_ghost(Monster *source) {
    Monster *copy = allocate_ghost();
    copy->health = source->health;
    copy->speed = source->speed;
    copy->clone = clone_ghost;
    return copy;
}
\end{lstlisting}

\Needspace{8\baselineskip}
\subsection{4.3 Function-pointer Spawner}

\textbf{Purpose:} Replace cloneable prototypes with first-class creation functions.

\begin{lstlisting}
typedef Monster *(*SpawnFn)(void);

typedef struct SpawnFunctionSpawner {
    SpawnFn spawn;
} SpawnFunctionSpawner;

Monster *spawn_from_function(SpawnFunctionSpawner *spawner) {
    return spawner->spawn();
}

Monster *spawn_ghost(void) {
    return allocate_ghost_with_defaults();
}
\end{lstlisting}

\Needspace{8\baselineskip}
\subsection{4.4 Data Prototype Delegation}

\textbf{Purpose:} Let data records inherit missing properties from a named prototype record.

\begin{lstlisting}
typedef struct DataRecord DataRecord;
struct DataRecord {
    const char *name;
    const char *prototype_name;
    Map properties;
};

DataRecord *find_record(const char *name);

Value lookup_property(DataRecord *record, const char *property_name) {
    if (map_contains(&record->properties, property_name)) {
        return map_get(&record->properties, property_name);
    }

    if (record->prototype_name != NULL) {
        DataRecord *prototype = find_record(record->prototype_name);
        if (prototype != NULL) {
            return lookup_property(prototype, property_name);
        }
    }

    return VALUE_MISSING;
}
\end{lstlisting}

\medskip\hrule\medskip

\clearpage
\section{5. Singleton and Single-instance Control Algorithms}

\Needspace{8\baselineskip}
\subsection{5.1 Lazy Single Instance}

\textbf{Purpose:} Create a service instance on first access.

\begin{lstlisting}
typedef struct FileSystem FileSystem;

FileSystem *file_system_instance(void) {
    static FileSystem *instance = NULL;

    if (instance == NULL) {
        instance = allocate_file_system();
    }

    return instance;
}
\end{lstlisting}

\Needspace{8\baselineskip}
\subsection{5.2 Runtime Single-instance Guard}

\textbf{Purpose:} Enforce one instance without granting global access.

\begin{lstlisting}
typedef struct SingleFileSystem {
    int internal_state;
} SingleFileSystem;

static Bool file_system_already_exists = FALSE;

void single_file_system_init(SingleFileSystem *fs) {
    assert(file_system_already_exists == FALSE);
    file_system_already_exists = TRUE;
    fs->internal_state = 0;
}

void single_file_system_destroy(SingleFileSystem *fs) {
    file_system_already_exists = FALSE;
}
\end{lstlisting}

\Needspace{8\baselineskip}
\subsection{5.3 Access Through Existing Global Context}

\textbf{Purpose:} Reduce singleton count by routing service access through a single application context.

\begin{lstlisting}
typedef struct GameContext {
    FileSystem *file_system;
    AudioService *audio;
    LogService *log;
} GameContext;

static GameContext *global_game;

AudioService *game_audio(void) {
    return global_game->audio;
}

void play_sound_from_context(int sound_id) {
    audio_play(game_audio(), sound_id);
}
\end{lstlisting}

\medskip\hrule\medskip

\clearpage
\section{6. State Pattern Algorithms}

\Needspace{8\baselineskip}
\subsection{6.1 Enum-based Finite State Machine}

\textbf{Purpose:} Replace invalid Boolean flag combinations with a single explicit state.

\begin{lstlisting}
typedef enum HeroineStateId {
    STATE_STANDING,
    STATE_JUMPING,
    STATE_DUCKING,
    STATE_DIVING
} HeroineStateId;

typedef struct Heroine {
    HeroineStateId state;
    float y_velocity;
    int charge_time;
} Heroine;

void heroine_handle_input_enum(Heroine *h, Input input) {
    switch (h->state) {
        case STATE_STANDING:
            if (input == PRESS_B) {
                h->state = STATE_JUMPING;
                h->y_velocity = JUMP_VELOCITY;
                heroine_set_graphics(h, IMAGE_JUMP);
            } else if (input == PRESS_DOWN) {
                h->state = STATE_DUCKING;
                h->charge_time = 0;
                heroine_set_graphics(h, IMAGE_DUCK);
            }
            break;

        case STATE_JUMPING:
            if (input == PRESS_DOWN) {
                h->state = STATE_DIVING;
                heroine_set_graphics(h, IMAGE_DIVE);
            }
            break;

        case STATE_DUCKING:
            if (input == RELEASE_DOWN) {
                h->state = STATE_STANDING;
                heroine_set_graphics(h, IMAGE_STAND);
            }
            break;

        case STATE_DIVING:
            break;
    }
}
\end{lstlisting}

\Needspace{8\baselineskip}
\subsection{6.2 State Object Dispatch}

\textbf{Purpose:} Encapsulate per-state behavior and state-specific data in separate objects.

\begin{lstlisting}
typedef struct HeroineState HeroineState;
struct HeroineState {
    HeroineState *(*handle_input)(HeroineState *self, Heroine *heroine, Input input);
    void (*update)(HeroineState *self, Heroine *heroine);
    void (*enter)(HeroineState *self, Heroine *heroine);
    void (*exit)(HeroineState *self, Heroine *heroine);
};

struct Heroine {
    HeroineState *state;
    float y_velocity;
};

void heroine_transition(Heroine *heroine, HeroineState *next) {
    if (next == NULL || next == heroine->state) return;

    heroine->state->exit(heroine->state, heroine);
    destroy_state_if_owned(heroine->state);
    heroine->state = next;
    heroine->state->enter(heroine->state, heroine);
}

void heroine_handle_input_state_object(Heroine *heroine, Input input) {
    HeroineState *next = heroine->state->handle_input(heroine->state, heroine, input);
    heroine_transition(heroine, next);
}

void heroine_update_state_object(Heroine *heroine) {
    heroine->state->update(heroine->state, heroine);
}
\end{lstlisting}

\Needspace{8\baselineskip}
\subsection{6.3 Ducking State with Local Timer}

\textbf{Purpose:} Keep charge-time state local to the ducking state.

\begin{lstlisting}
typedef struct DuckingState {
    HeroineState base;
    int charge_time;
} DuckingState;

HeroineState *ducking_handle_input(HeroineState *base, Heroine *heroine, Input input) {
    if (input == RELEASE_DOWN) {
        return make_standing_state();
    }
    return NULL;
}

void ducking_update(HeroineState *base, Heroine *heroine) {
    DuckingState *state = (DuckingState *)base;
    state->charge_time++;
    if (state->charge_time > MAX_CHARGE) {
        heroine_super_bomb(heroine);
    }
}

void ducking_enter(HeroineState *base, Heroine *heroine) {
    DuckingState *state = (DuckingState *)base;
    state->charge_time = 0;
    heroine_set_graphics(heroine, IMAGE_DUCK);
}
\end{lstlisting}

\Needspace{8\baselineskip}
\subsection{6.4 Concurrent State Machines}

\textbf{Purpose:} Avoid combinatorial explosion by modeling independent state dimensions separately.

\begin{lstlisting}
typedef struct HeroineConcurrent {
    HeroineState *movement_state;
    HeroineState *equipment_state;
} HeroineConcurrent;

void heroine_handle_concurrent(HeroineConcurrent *h, Input input) {
    HeroineState *next_movement = h->movement_state->handle_input(h->movement_state, (Heroine *)h, input);
    HeroineState *next_equipment = h->equipment_state->handle_input(h->equipment_state, (Heroine *)h, input);

    if (next_movement != NULL) h->movement_state = next_movement;
    if (next_equipment != NULL) h->equipment_state = next_equipment;
}
\end{lstlisting}

\Needspace{8\baselineskip}
\subsection{6.5 Hierarchical State Machine}

\textbf{Purpose:} Let substates defer unhandled input to superstates.

\begin{lstlisting}
struct HeroineState {
    HeroineState *superstate;
    Bool (*handle_input)(HeroineState *self, Heroine *heroine, Input input);
};

void hsm_handle_input(Heroine *heroine, Input input) {
    HeroineState *state = heroine->state;

    while (state != NULL) {
        if (state->handle_input(state, heroine, input)) {
            return;
        }
        state = state->superstate;
    }
}
\end{lstlisting}

\Needspace{8\baselineskip}
\subsection{6.6 Pushdown Automaton}

\textbf{Purpose:} Return automatically to a previous state after a temporary state finishes.

\begin{lstlisting}
#define MAX_STATE_STACK 32

typedef struct StateStack {
    HeroineState *items[MAX_STATE_STACK];
    int count;
} StateStack;

HeroineState *current_state(StateStack *stack) {
    return stack->items[stack->count - 1];
}

void push_state(StateStack *stack, HeroineState *state) {
    stack->items[stack->count++] = state;
    state->enter(state, NULL);
}

void pop_state(StateStack *stack) {
    if (stack->count == 0) return;
    HeroineState *old = stack->items[--stack->count];
    old->exit(old, NULL);
}

void fire_weapon_with_return(StateStack *stack) {
    push_state(stack, make_firing_state());
}

void firing_animation_finished(StateStack *stack) {
    pop_state(stack); /* resumes the state beneath firing */
}
\end{lstlisting}

\medskip\hrule\medskip

\clearpage
\section{7. Double Buffer Algorithms}

\Needspace{8\baselineskip}
\subsection{7.1 Framebuffer Double Buffer}

\textbf{Purpose:} Render a complete next frame while the display reads the current frame.

\begin{lstlisting}
#define WIDTH 160
#define HEIGHT 120

typedef struct Framebuffer {
    char pixels[WIDTH * HEIGHT];
} Framebuffer;

typedef struct SceneBuffers {
    Framebuffer buffers[2];
    Framebuffer *current;
    Framebuffer *next;
} SceneBuffers;

void framebuffer_clear(Framebuffer *buffer, char color) {
    for (int i = 0; i < WIDTH * HEIGHT; i++) {
        buffer->pixels[i] = color;
    }
}

void framebuffer_draw(Framebuffer *buffer, int x, int y, char color) {
    buffer->pixels[WIDTH * y + x] = color;
}

void scene_buffers_init(SceneBuffers *scene) {
    scene->current = &scene->buffers[0];
    scene->next = &scene->buffers[1];
}

void scene_swap(SceneBuffers *scene) {
    Framebuffer *temp = scene->current;
    scene->current = scene->next;
    scene->next = temp;
}

void scene_draw(SceneBuffers *scene) {
    framebuffer_clear(scene->next, WHITE);
    draw_all_visible_objects(scene->next);
    scene_swap(scene);
}

Framebuffer *scene_get_display_buffer(SceneBuffers *scene) {
    return scene->current;
}
\end{lstlisting}

\Needspace{8\baselineskip}
\subsection{7.2 Fine-grained Actor Double Buffer}

\textbf{Purpose:} Make same-frame actor updates appear simultaneous by reading current flags and writing next flags.

\begin{lstlisting}
typedef struct SlapActor {
    Bool current_slapped;
    Bool next_slapped;
    void (*update)(struct SlapActor *self);
} SlapActor;

void actor_slap(SlapActor *actor) {
    actor->next_slapped = TRUE;
}

Bool actor_was_slapped(SlapActor *actor) {
    return actor->current_slapped;
}

void actor_swap(SlapActor *actor) {
    actor->current_slapped = actor->next_slapped;
    actor->next_slapped = FALSE;
}

void stage_update(SlapActor **actors, int count) {
    for (int i = 0; i < count; i++) {
        actors[i]->update(actors[i]);
    }

    for (int i = 0; i < count; i++) {
        actor_swap(actors[i]);
    }
}
\end{lstlisting}

\medskip\hrule\medskip

\clearpage
\section{8. Game Loop Algorithms}

\Needspace{8\baselineskip}
\subsection{8.1 Minimal Game Loop}

\textbf{Purpose:} Keep the game running until the player exits.

\begin{lstlisting}
void run_game_simple(void) {
    while (!game_should_quit()) {
        process_input();
        update_game();
        render_game();
    }
}
\end{lstlisting}

\Needspace{8\baselineskip}
\subsection{8.2 Variable Timestep Loop}

\textbf{Purpose:} Scale simulation progress by elapsed real time.

\begin{lstlisting}
void run_game_variable_timestep(void) {
    double previous = now_seconds();

    while (!game_should_quit()) {
        double current = now_seconds();
        double elapsed = current - previous;
        previous = current;

        process_input();
        update_game(elapsed);
        render_game();
    }
}
\end{lstlisting}

\Needspace{8\baselineskip}
\subsection{8.3 Fixed Timestep Loop with Sleep}

\textbf{Purpose:} Keep gameplay deterministic by advancing in fixed increments.

\begin{lstlisting}
#define FRAME_SECONDS (1.0 / 60.0)

void run_game_fixed_timestep(void) {
    double previous = now_seconds();
    double lag = 0.0;

    while (!game_should_quit()) {
        double current = now_seconds();
        double elapsed = current - previous;
        previous = current;
        lag += elapsed;

        process_input();

        while (lag >= FRAME_SECONDS) {
            update_game(FRAME_SECONDS);
            lag -= FRAME_SECONDS;
        }

        render_game();
    }
}
\end{lstlisting}

\Needspace{8\baselineskip}
\subsection{8.4 Fixed Timestep with Render Interpolation}

\textbf{Purpose:} Render smoothly between discrete simulation ticks.

\begin{lstlisting}
void run_game_fixed_with_interpolation(void) {
    double previous = now_seconds();
    double lag = 0.0;

    while (!game_should_quit()) {
        double current = now_seconds();
        double elapsed = current - previous;
        previous = current;
        lag += elapsed;

        process_input();

        while (lag >= FRAME_SECONDS) {
            store_previous_simulation_state();
            update_game(FRAME_SECONDS);
            lag -= FRAME_SECONDS;
        }

        double alpha = lag / FRAME_SECONDS;
        render_interpolated(alpha);
    }
}
\end{lstlisting}

\medskip\hrule\medskip

\clearpage
\section{9. Update Method Algorithms}

\Needspace{8\baselineskip}
\subsection{9.1 Per-object Update Loop}

\textbf{Purpose:} Simulate independent entities one frame at a time.

\begin{lstlisting}
typedef struct GameObject GameObject;
struct GameObject {
    void (*update)(GameObject *self);
    Bool active;
};

void update_objects(GameObject **objects, int count) {
    for (int i = 0; i < count; i++) {
        if (objects[i]->active) {
            objects[i]->update(objects[i]);
        }
    }
}
\end{lstlisting}

\Needspace{8\baselineskip}
\subsection{9.2 Patrol Skeleton Update}

\textbf{Purpose:} Encode object-local temporal behavior as fields plus an update method.

\begin{lstlisting}
typedef struct Skeleton {
    GameObject base;
    Vec2 position;
    int direction; /* -1 or +1 */
} Skeleton;

void skeleton_update(GameObject *base) {
    Skeleton *s = (Skeleton *)base;
    s->position.x += s->direction * SKELETON_SPEED;

    if (s->position.x < LEFT_LIMIT) {
        s->position.x = LEFT_LIMIT;
        s->direction = +1;
    } else if (s->position.x > RIGHT_LIMIT) {
        s->position.x = RIGHT_LIMIT;
        s->direction = -1;
    }
}
\end{lstlisting}

\Needspace{8\baselineskip}
\subsection{9.3 Deferred Add/Remove During Update}

\textbf{Purpose:} Avoid corrupting an update loop when objects are created or destroyed mid-iteration.

\begin{lstlisting}
typedef struct ObjectWorld {
    GameObject *objects[4096];
    int count;
    GameObject *pending_add[1024];
    int pending_add_count;
    GameObject *pending_remove[1024];
    int pending_remove_count;
    Bool updating;
} ObjectWorld;

void world_add(ObjectWorld *world, GameObject *object) {
    if (world->updating) {
        world->pending_add[world->pending_add_count++] = object;
    } else {
        world->objects[world->count++] = object;
    }
}

void world_remove(ObjectWorld *world, GameObject *object) {
    if (world->updating) {
        world->pending_remove[world->pending_remove_count++] = object;
    } else {
        remove_object_immediately(world, object);
    }
}

void world_update(ObjectWorld *world) {
    world->updating = TRUE;
    for (int i = 0; i < world->count; i++) {
        if (world->objects[i]->active) world->objects[i]->update(world->objects[i]);
    }
    world->updating = FALSE;

    for (int i = 0; i < world->pending_remove_count; i++) {
        remove_object_immediately(world, world->pending_remove[i]);
    }
    world->pending_remove_count = 0;

    for (int i = 0; i < world->pending_add_count; i++) {
        world->objects[world->count++] = world->pending_add[i];
    }
    world->pending_add_count = 0;
}
\end{lstlisting}

\medskip\hrule\medskip

\clearpage
\section{10. Bytecode Algorithms}

\Needspace{8\baselineskip}
\subsection{10.1 Stack Virtual Machine}

\textbf{Purpose:} Interpret compact bytecode instructions with an operand stack.

\begin{lstlisting}
typedef enum OpCode {
    OP_LITERAL,
    OP_GET_HEALTH,
    OP_SET_HEALTH,
    OP_GET_WISDOM,
    OP_SET_WISDOM,
    OP_ADD,
    OP_DIVIDE,
    OP_PLAY_SOUND,
    OP_SPAWN_PARTICLES,
    OP_RETURN
} OpCode;

typedef struct VM {
    unsigned char *code;
    int ip;
    int stack[256];
    int stack_top;
} VM;

void vm_push(VM *vm, int value) {
    vm->stack[vm->stack_top++] = value;
}

int vm_pop(VM *vm) {
    return vm->stack[--vm->stack_top];
}

unsigned char vm_read_byte(VM *vm) {
    return vm->code[vm->ip++];
}

void vm_execute(VM *vm) {
    for (;;) {
        OpCode op = (OpCode)vm_read_byte(vm);

        switch (op) {
            case OP_LITERAL: {
                int value = vm_read_byte(vm);
                vm_push(vm, value);
                break;
            }

            case OP_GET_HEALTH: {
                int wizard = vm_pop(vm);
                vm_push(vm, wizard_health(wizard));
                break;
            }

            case OP_SET_HEALTH: {
                int amount = vm_pop(vm);
                int wizard = vm_pop(vm);
                wizard_set_health(wizard, amount);
                break;
            }

            case OP_GET_WISDOM: {
                int wizard = vm_pop(vm);
                vm_push(vm, wizard_wisdom(wizard));
                break;
            }

            case OP_SET_WISDOM: {
                int amount = vm_pop(vm);
                int wizard = vm_pop(vm);
                wizard_set_wisdom(wizard, amount);
                break;
            }

            case OP_ADD: {
                int b = vm_pop(vm);
                int a = vm_pop(vm);
                vm_push(vm, a + b);
                break;
            }

            case OP_DIVIDE: {
                int divisor = vm_pop(vm);
                int dividend = vm_pop(vm);
                if (divisor == 0) vm_push(vm, 0);
                else vm_push(vm, dividend / divisor);
                break;
            }

            case OP_PLAY_SOUND: {
                int sound_id = vm_pop(vm);
                play_sound(sound_id);
                break;
            }

            case OP_SPAWN_PARTICLES: {
                int particle_id = vm_pop(vm);
                spawn_particles(particle_id);
                break;
            }

            case OP_RETURN:
                return;
        }
    }
}
\end{lstlisting}

\Needspace{8\baselineskip}
\subsection{10.2 Bytecode Disassembly}

\textbf{Purpose:} Turn bytecode into readable diagnostics.

\begin{lstlisting}
void disassemble(unsigned char *code, int length) {
    int ip = 0;
    while (ip < length) {
        int offset = ip;
        OpCode op = (OpCode)code[ip++];

        print_int(offset);
        print_string(": ");

        switch (op) {
            case OP_LITERAL:
                print_string("LITERAL ");
                print_int(code[ip++]);
                break;
            case OP_ADD:
                print_string("ADD");
                break;
            case OP_DIVIDE:
                print_string("DIVIDE");
                break;
            case OP_RETURN:
                print_string("RETURN");
                break;
            default:
                print_string("<unknown>");
                break;
        }
        print_newline();
    }
}
\end{lstlisting}

\medskip\hrule\medskip

\clearpage
\section{11. Subclass Sandbox Algorithms}

\Needspace{8\baselineskip}
\subsection{11.1 Base-class Sandbox Interface}

\textbf{Purpose:} Let derived gameplay classes implement behavior through a safe restricted API.

\begin{lstlisting}
typedef struct Superpower Superpower;
struct Superpower {
    void (*activate)(Superpower *self);
};

void power_move(float x, float y, float z) {
    engine_move_actor(current_actor(), x, y, z);
}

void power_play_sound(int sound_id) {
    audio_queue_sound(sound_id);
}

void power_spawn_particles(int particle_id, Vec3 position) {
    particle_system_spawn(particle_id, position);
}
\end{lstlisting}

\Needspace{8\baselineskip}
\subsection{11.2 Derived Power Behavior}

\textbf{Purpose:} Compose high-level primitive operations rather than reaching into engine subsystems directly.

\begin{lstlisting}
typedef struct SkyLaunchPower {
    Superpower base;
    float impulse;
} SkyLaunchPower;

void sky_launch_activate(Superpower *base) {
    SkyLaunchPower *power = (SkyLaunchPower *)base;
    power_play_sound(SOUND_LAUNCH);
    power_spawn_particles(PARTICLE_DUST_RING, actor_position(current_actor()));
    power_move(0.0f, power->impulse, 0.0f);
}
\end{lstlisting}

\Needspace{8\baselineskip}
\subsection{11.3 Sandbox Operation with Shared Service}

\textbf{Purpose:} Centralize subsystem access and policy in the base layer.

\begin{lstlisting}
static AudioService *sandbox_audio;

void sandbox_set_audio(AudioService *audio) {
    sandbox_audio = audio;
}

void power_play_sound_restricted(int sound_id) {
    if (sandbox_audio == NULL) return;
    if (!sound_is_allowed_for_power(sound_id)) return;
    audio_play(sandbox_audio, sound_id);
}
\end{lstlisting}

\medskip\hrule\medskip

\clearpage
\section{12. Type Object Algorithms}

\Needspace{8\baselineskip}
\subsection{12.1 Monster Delegates Behavior to Breed}

\textbf{Purpose:} Move type-specific attributes and methods into runtime type objects.

\begin{lstlisting}
typedef struct Breed Breed;

typedef struct MonsterInstance {
    Breed *breed;
    int health;
} MonsterInstance;

struct Breed {
    const char *name;
    int starting_health;
    const char *attack_text;
    Breed *parent;
};

void monster_init(MonsterInstance *monster, Breed *breed) {
    monster->breed = breed;
    monster->health = breed->starting_health;
}

void monster_attack(MonsterInstance *monster) {
    print_string(resolve_attack_text(monster->breed));
}
\end{lstlisting}

\Needspace{8\baselineskip}
\subsection{12.2 Breed Property Inheritance}

\textbf{Purpose:} Resolve missing breed fields through a parent chain.

\begin{lstlisting}
const char *resolve_attack_text(Breed *breed) {
    Breed *current = breed;
    while (current != NULL) {
        if (current->attack_text != NULL) {
            return current->attack_text;
        }
        current = current->parent;
    }
    return "attacks";
}

int resolve_starting_health(Breed *breed) {
    Breed *current = breed;
    while (current != NULL) {
        if (current->starting_health >= 0) {
            return current->starting_health;
        }
        current = current->parent;
    }
    return DEFAULT_MONSTER_HEALTH;
}
\end{lstlisting}

\Needspace{8\baselineskip}
\subsection{12.3 Breed-based Spawn}

\textbf{Purpose:} Let type objects manufacture configured instances.

\begin{lstlisting}
MonsterInstance *breed_new_monster(Breed *breed) {
    MonsterInstance *monster = allocate_monster();
    monster->breed = breed;
    monster->health = resolve_starting_health(breed);
    return monster;
}
\end{lstlisting}

\medskip\hrule\medskip

\clearpage
\section{13. Component Pattern Algorithms}

\Needspace{8\baselineskip}
\subsection{13.1 Entity Delegates to Components}

\textbf{Purpose:} Compose behavior using independent component objects.

\begin{lstlisting}
typedef struct InputComponent InputComponent;
typedef struct PhysicsComponent PhysicsComponent;
typedef struct GraphicsComponent GraphicsComponent;

typedef struct Entity {
    InputComponent *input;
    PhysicsComponent *physics;
    GraphicsComponent *graphics;
} Entity;

void entity_update(Entity *entity, World *world) {
    input_component_update(entity->input, entity);
    physics_component_update(entity->physics, entity, world);
    graphics_component_update(entity->graphics, entity);
}
\end{lstlisting}

\Needspace{8\baselineskip}
\subsection{13.2 Input Component Produces Intent}

\textbf{Purpose:} Isolate input mapping from physics and rendering.

\begin{lstlisting}
typedef struct Intent {
    float move_x;
    float move_y;
    Bool jump;
} Intent;

typedef struct InputComponent {
    Intent intent;
    void (*update)(struct InputComponent *self, Entity *entity);
} InputComponent;

void player_input_update(InputComponent *input, Entity *entity) {
    input->intent.move_x = axis_value(AXIS_HORIZONTAL);
    input->intent.move_y = axis_value(AXIS_VERTICAL);
    input->intent.jump = is_pressed(BUTTON_JUMP);
}

void ai_input_update(InputComponent *input, Entity *entity) {
    input->intent = ai_decide_intent(entity);
}
\end{lstlisting}

\Needspace{8\baselineskip}
\subsection{13.3 Physics Reads Intent and Updates Transform}

\textbf{Purpose:} Let physics remain independent of whether input comes from a player or AI.

\begin{lstlisting}
void physics_component_update(PhysicsComponent *physics, Entity *entity, World *world) {
    Intent intent = entity->input->intent;
    physics->velocity.x += intent.move_x * physics->acceleration;
    physics->velocity.y += intent.move_y * physics->acceleration;

    if (intent.jump && physics->on_ground) {
        physics->velocity.y = JUMP_VELOCITY;
    }

    physics_integrate(physics, entity, world);
}
\end{lstlisting}

\Needspace{8\baselineskip}
\subsection{13.4 Component Messaging}

\textbf{Purpose:} Send local events among components without creating direct dependencies between all components.

\begin{lstlisting}
typedef struct ComponentMessage {
    int type;
    int int_arg;
    Vec3 vec_arg;
} ComponentMessage;

void entity_send(Entity *entity, ComponentMessage message) {
    input_component_receive(entity->input, entity, message);
    physics_component_receive(entity->physics, entity, message);
    graphics_component_receive(entity->graphics, entity, message);
}
\end{lstlisting}

\medskip\hrule\medskip

\clearpage
\section{14. Event Queue Algorithms}

\Needspace{8\baselineskip}
\subsection{14.1 Ring-buffer Event Queue}

\textbf{Purpose:} Decouple event producers from event consumers in time.

\begin{lstlisting}
#define EVENT_QUEUE_SIZE 256

typedef struct Event {
    EventType type;
    EntityId source;
    EntityId target;
    int payload0;
    int payload1;
} Event;

typedef struct EventQueue {
    Event events[EVENT_QUEUE_SIZE];
    int head;
    int tail;
} EventQueue;

Bool event_queue_is_empty(EventQueue *q) {
    return q->head == q->tail;
}

Bool event_queue_is_full(EventQueue *q) {
    return ((q->tail + 1) % EVENT_QUEUE_SIZE) == q->head;
}

void event_enqueue(EventQueue *q, Event event) {
    if (event_queue_is_full(q)) {
        drop_or_expand_event_queue(q, event);
        return;
    }
    q->events[q->tail] = event;
    q->tail = (q->tail + 1) % EVENT_QUEUE_SIZE;
}

Bool event_dequeue(EventQueue *q, Event *out) {
    if (event_queue_is_empty(q)) return FALSE;
    *out = q->events[q->head];
    q->head = (q->head + 1) % EVENT_QUEUE_SIZE;
    return TRUE;
}
\end{lstlisting}

\Needspace{8\baselineskip}
\subsection{14.2 Process Events Later}

\textbf{Purpose:} Drain the event queue at a controlled point in the game loop.

\begin{lstlisting}
void process_event_queue(EventQueue *q) {
    Event event;
    while (event_dequeue(q, &event)) {
        switch (event.type) {
            case EVENT_ENEMY_DIED:
                tutorial_on_enemy_died(event.source);
                score_on_enemy_died(event.source);
                break;

            case EVENT_SOUND_REQUESTED:
                audio_process_request(event.payload0, event.payload1);
                break;
        }
    }
}
\end{lstlisting}

\Needspace{8\baselineskip}
\subsection{14.3 Audio Request Queue with Duplicate Coalescing}

\textbf{Purpose:} Avoid playing the same sound too many times in one frame.

\begin{lstlisting}
#define MAX_PENDING_SOUNDS 64

typedef struct SoundRequest {
    int sound_id;
    int volume;
} SoundRequest;

typedef struct AudioQueue {
    SoundRequest pending[MAX_PENDING_SOUNDS];
    int count;
} AudioQueue;

void audio_enqueue(AudioQueue *q, int sound_id, int volume) {
    for (int i = 0; i < q->count; i++) {
        if (q->pending[i].sound_id == sound_id) {
            if (volume > q->pending[i].volume) {
                q->pending[i].volume = volume;
            }
            return;
        }
    }

    if (q->count < MAX_PENDING_SOUNDS) {
        q->pending[q->count++] = (SoundRequest){ sound_id, volume };
    }
}

void audio_process(AudioQueue *q) {
    for (int i = 0; i < q->count; i++) {
        audio_play(q->pending[i].sound_id, q->pending[i].volume);
    }
    q->count = 0;
}
\end{lstlisting}

\medskip\hrule\medskip

\clearpage
\section{15. Service Locator Algorithms}

\Needspace{8\baselineskip}
\subsection{15.1 Locate and Provide Service}

\textbf{Purpose:} Make a service globally reachable through an abstract interface.

\begin{lstlisting}
typedef struct AudioService AudioService;
struct AudioService {
    void (*play_sound)(AudioService *self, int sound_id);
    void (*stop_sound)(AudioService *self, int sound_id);
};

static AudioService *located_audio = NULL;

void service_provide_audio(AudioService *service) {
    located_audio = service;
}

AudioService *service_audio(void) {
    return located_audio;
}
\end{lstlisting}

\Needspace{8\baselineskip}
\subsection{15.2 Null Service}

\textbf{Purpose:} Guarantee callers always receive a safe object.

\begin{lstlisting}
void null_audio_play(AudioService *self, int sound_id) {
    /* intentionally silent */
}

void null_audio_stop(AudioService *self, int sound_id) {
    /* intentionally silent */
}

AudioService *null_audio_service(void) {
    static AudioService null_service = { null_audio_play, null_audio_stop };
    return &null_service;
}

AudioService *service_audio_safe(void) {
    if (located_audio == NULL) return null_audio_service();
    return located_audio;
}
\end{lstlisting}

\Needspace{8\baselineskip}
\subsection{15.3 Logged Service Decorator}

\textbf{Purpose:} Wrap a service with diagnostics while preserving the service interface.

\begin{lstlisting}
typedef struct LoggedAudioService {
    AudioService base;
    AudioService *wrapped;
} LoggedAudioService;

void logged_audio_play(AudioService *base, int sound_id) {
    LoggedAudioService *self = (LoggedAudioService *)base;
    log_int("play sound", sound_id);
    self->wrapped->play_sound(self->wrapped, sound_id);
}

void logged_audio_stop(AudioService *base, int sound_id) {
    LoggedAudioService *self = (LoggedAudioService *)base;
    log_int("stop sound", sound_id);
    self->wrapped->stop_sound(self->wrapped, sound_id);
}
\end{lstlisting}

\medskip\hrule\medskip

\clearpage
\section{16. Data Locality Algorithms}

\Needspace{8\baselineskip}
\subsection{16.1 Packed Component Update}

\textbf{Purpose:} Iterate contiguous component data instead of pointer-chasing scattered objects.

\begin{lstlisting}
typedef struct Particle {
    Vec3 position;
    Vec3 velocity;
    float lifetime;
    Bool active;
} Particle;

typedef struct ParticleSystem {
    Particle particles[10000];
    int active_count;
} ParticleSystem;

void update_particles_packed(ParticleSystem *system, float dt) {
    for (int i = 0; i < system->active_count; i++) {
        Particle *p = &system->particles[i];
        p->position.x += p->velocity.x * dt;
        p->position.y += p->velocity.y * dt;
        p->position.z += p->velocity.z * dt;
        p->lifetime -= dt;
    }
}
\end{lstlisting}

\Needspace{8\baselineskip}
\subsection{16.2 Active-prefix Partition}

\textbf{Purpose:} Keep active objects contiguous by swapping activated/deactivated objects across a boundary.

\begin{lstlisting}
void activate_particle(ParticleSystem *system, int index) {
    if (index < system->active_count) return;

    swap_particles(&system->particles[index], &system->particles[system->active_count]);
    system->active_count++;
}

void deactivate_particle(ParticleSystem *system, int index) {
    if (index >= system->active_count) return;

    system->active_count--;
    swap_particles(&system->particles[index], &system->particles[system->active_count]);
}
\end{lstlisting}

\Needspace{8\baselineskip}
\subsection{16.3 Hot/Cold Split}

\textbf{Purpose:} Keep frequently accessed fields in a compact hot array and rarely used fields elsewhere.

\begin{lstlisting}
typedef struct AiHot {
    Vec3 position;
    Vec3 velocity;
    int state;
    int cold_index;
} AiHot;

typedef struct AiCold {
    char debug_name[64];
    char script_name[64];
    Inventory inventory;
} AiCold;

typedef struct AiStorage {
    AiHot hot[4096];
    AiCold cold[4096];
    int count;
} AiStorage;

void update_ai_hot(AiStorage *storage, float dt) {
    for (int i = 0; i < storage->count; i++) {
        ai_integrate(&storage->hot[i], dt);
    }
}
\end{lstlisting}

\medskip\hrule\medskip

\clearpage
\section{17. Dirty Flag Algorithms}

\Needspace{8\baselineskip}
\subsection{17.1 Dirty Transform Node}

\textbf{Purpose:} Cache world transforms and recompute only after local or ancestor transform changes.

\begin{lstlisting}
typedef struct Transform Transform;

typedef struct SceneNode SceneNode;
struct SceneNode {
    Transform local;
    Transform world;
    Bool dirty;
    SceneNode *parent;
    SceneNode *children[64];
    int child_count;
};

void node_mark_dirty(SceneNode *node) {
    if (node->dirty) return;

    node->dirty = TRUE;
    for (int i = 0; i < node->child_count; i++) {
        node_mark_dirty(node->children[i]);
    }
}

void node_set_local(SceneNode *node, Transform local) {
    node->local = local;
    node_mark_dirty(node);
}
\end{lstlisting}

\Needspace{8\baselineskip}
\subsection{17.2 Lazy World Transform Resolution}

\textbf{Purpose:} Resolve derived data only when requested.

\begin{lstlisting}
Transform node_world_transform(SceneNode *node) {
    if (node->dirty) {
        if (node->parent == NULL) {
            node->world = node->local;
        } else {
            Transform parent_world = node_world_transform(node->parent);
            node->world = transform_multiply(parent_world, node->local);
        }
        node->dirty = FALSE;
    }

    return node->world;
}
\end{lstlisting}

\Needspace{8\baselineskip}
\subsection{17.3 Eager Dirty Flush}

\textbf{Purpose:} Recompute dirty transforms in one traversal at a controlled point.

\begin{lstlisting}
void flush_world_transforms(SceneNode *node, Transform parent_world, Bool parent_dirty) {
    Bool must_recompute = parent_dirty || node->dirty;

    if (must_recompute) {
        node->world = transform_multiply(parent_world, node->local);
        node->dirty = FALSE;
    }

    for (int i = 0; i < node->child_count; i++) {
        flush_world_transforms(node->children[i], node->world, must_recompute);
    }
}
\end{lstlisting}

\medskip\hrule\medskip

\clearpage
\section{18. Object Pool Algorithms}

\Needspace{8\baselineskip}
\subsection{18.1 Fixed-size Object Pool with Free List}

\textbf{Purpose:} Allocate and free objects from a preallocated fixed array.

\begin{lstlisting}
#define POOL_SIZE 1024

typedef struct PoolObject PoolObject;
struct PoolObject {
    Bool in_use;
    union {
        struct {
            Vec3 position;
            Vec3 velocity;
            float lifetime;
        } live;
        PoolObject *next_free;
    } data;
};

typedef struct ObjectPool {
    PoolObject objects[POOL_SIZE];
    PoolObject *free_head;
} ObjectPool;

void pool_init(ObjectPool *pool) {
    pool->free_head = &pool->objects[0];

    for (int i = 0; i < POOL_SIZE - 1; i++) {
        pool->objects[i].in_use = FALSE;
        pool->objects[i].data.next_free = &pool->objects[i + 1];
    }

    pool->objects[POOL_SIZE - 1].in_use = FALSE;
    pool->objects[POOL_SIZE - 1].data.next_free = NULL;
}

PoolObject *pool_create(ObjectPool *pool) {
    if (pool->free_head == NULL) return NULL;

    PoolObject *object = pool->free_head;
    pool->free_head = object->data.next_free;
    object->in_use = TRUE;
    return object;
}

void pool_destroy(ObjectPool *pool, PoolObject *object) {
    object->in_use = FALSE;
    object->data.next_free = pool->free_head;
    pool->free_head = object;
}
\end{lstlisting}

\Needspace{8\baselineskip}
\subsection{18.2 Pool Update}

\textbf{Purpose:} Update only live pooled objects and return expired ones to the pool.

\begin{lstlisting}
void pool_update(ObjectPool *pool, float dt) {
    for (int i = 0; i < POOL_SIZE; i++) {
        PoolObject *object = &pool->objects[i];
        if (!object->in_use) continue;

        object->data.live.position.x += object->data.live.velocity.x * dt;
        object->data.live.position.y += object->data.live.velocity.y * dt;
        object->data.live.position.z += object->data.live.velocity.z * dt;
        object->data.live.lifetime -= dt;

        if (object->data.live.lifetime <= 0.0f) {
            pool_destroy(pool, object);
        }
    }
}
\end{lstlisting}

\Needspace{8\baselineskip}
\subsection{18.3 Particle Spawn from Pool}

\textbf{Purpose:} Convert a free pool slot into a live particle with initialized state.

\begin{lstlisting}
PoolObject *spawn_particle(ObjectPool *pool, Vec3 position, Vec3 velocity, float lifetime) {
    PoolObject *particle = pool_create(pool);
    if (particle == NULL) return NULL;

    particle->data.live.position = position;
    particle->data.live.velocity = velocity;
    particle->data.live.lifetime = lifetime;
    return particle;
}
\end{lstlisting}

\medskip\hrule\medskip

\clearpage
\section{19. Spatial Partition Algorithms}

\Needspace{8\baselineskip}
\subsection{19.1 Uniform Grid with Intrusive Cell Lists}

\textbf{Purpose:} Restrict collision checks to nearby units by storing units in grid cells.

\begin{lstlisting}
#define GRID_W 128
#define GRID_H 128
#define CELL_SIZE 4.0f

typedef struct Unit Unit;
struct Unit {
    float x;
    float y;
    Unit *prev;
    Unit *next;
    int cell_x;
    int cell_y;
};

typedef struct Grid {
    Unit *cells[GRID_W][GRID_H];
} Grid;

int grid_cell_coord(float value) {
    return clamp_int((int)(value / CELL_SIZE), 0, GRID_W - 1);
}
\end{lstlisting}

\Needspace{8\baselineskip}
\subsection{19.2 Insert Unit into Grid}

\textbf{Purpose:} Link a unit into the cell corresponding to its position.

\begin{lstlisting}
void grid_insert(Grid *grid, Unit *unit) {
    int cell_x = grid_cell_coord(unit->x);
    int cell_y = grid_cell_coord(unit->y);

    unit->cell_x = cell_x;
    unit->cell_y = cell_y;
    unit->prev = NULL;
    unit->next = grid->cells[cell_x][cell_y];

    if (unit->next != NULL) {
        unit->next->prev = unit;
    }

    grid->cells[cell_x][cell_y] = unit;
}
\end{lstlisting}

\Needspace{8\baselineskip}
\subsection{19.3 Remove Unit from Grid}

\textbf{Purpose:} Unlink a unit from its current cell in constant time.

\begin{lstlisting}
void grid_remove(Grid *grid, Unit *unit) {
    int cell_x = unit->cell_x;
    int cell_y = unit->cell_y;

    if (unit->prev != NULL) {
        unit->prev->next = unit->next;
    } else {
        grid->cells[cell_x][cell_y] = unit->next;
    }

    if (unit->next != NULL) {
        unit->next->prev = unit->prev;
    }

    unit->prev = NULL;
    unit->next = NULL;
}
\end{lstlisting}

\Needspace{8\baselineskip}
\subsection{19.4 Move Unit Between Cells}

\textbf{Purpose:} Update grid membership only when a unit crosses a cell boundary.

\begin{lstlisting}
void grid_move(Grid *grid, Unit *unit, float x, float y) {
    int old_x = unit->cell_x;
    int old_y = unit->cell_y;
    int new_x = grid_cell_coord(x);
    int new_y = grid_cell_coord(y);

    unit->x = x;
    unit->y = y;

    if (old_x == new_x && old_y == new_y) return;

    grid_remove(grid, unit);
    grid_insert(grid, unit);
}
\end{lstlisting}

\Needspace{8\baselineskip}
\subsection{19.5 Handle Nearby Interactions}

\textbf{Purpose:} Test a unit only against units in its own and neighboring cells.

\begin{lstlisting}
void handle_unit_against_cell(Unit *unit, Unit *other) {
    while (other != NULL) {
        if (other != unit && units_overlap(unit, other)) {
            resolve_unit_interaction(unit, other);
        }
        other = other->next;
    }
}

void grid_handle_nearby(Grid *grid, Unit *unit) {
    int cx = unit->cell_x;
    int cy = unit->cell_y;

    for (int dx = -1; dx <= 1; dx++) {
        for (int dy = -1; dy <= 1; dy++) {
            int x = cx + dx;
            int y = cy + dy;
            if (x < 0 || y < 0 || x >= GRID_W || y >= GRID_H) continue;
            handle_unit_against_cell(unit, grid->cells[x][y]);
        }
    }
}
\end{lstlisting}

\Needspace{8\baselineskip}
\subsection{19.6 All-pairs Spatial Partition Pass}

\textbf{Purpose:} Process each unit while avoiding global O(n\textasciicircum{}2) pair checks.

\begin{lstlisting}
void grid_process_all_units(Grid *grid) {
    for (int x = 0; x < GRID_W; x++) {
        for (int y = 0; y < GRID_H; y++) {
            Unit *unit = grid->cells[x][y];
            while (unit != NULL) {
                grid_handle_nearby(grid, unit);
                unit = unit->next;
            }
        }
    }
}
\end{lstlisting}

\medskip\hrule\medskip

\clearpage
\section{20. Integrated Game-loop Skeleton Using the Extracted Patterns}

This final sketch shows how several algorithms fit together in a single runtime frame.

\begin{lstlisting}
void game_frame(GameContext *game, double dt) {
    /* 1. Poll platform input and turn it into commands. */
    ActorCommand *player_command = poll_actor_command(&game->input_bindings);
    if (player_command != NULL) {
        enqueue_command(&game->command_queue, player_command);
    }

    /* 2. Produce AI commands. */
    for (int i = 0; i < game->ai_actor_count; i++) {
        ActorCommand *ai_command = ai_choose_command(game->ai_actors[i]);
        if (ai_command != NULL) enqueue_command(&game->command_queue, ai_command);
    }

    /* 3. Execute commands against actors. */
    process_command_stream(&game->command_queue, game->player_actor);

    /* 4. Update game objects/components. */
    update_objects(game->objects, game->object_count);
    update_particles_packed(&game->particles, (float)dt);

    /* 5. Recompute dirty derived state at a controlled boundary. */
    flush_world_transforms(game->root_node, transform_identity(), TRUE);

    /* 6. Run spatial interaction queries. */
    grid_process_all_units(&game->spatial_grid);

    /* 7. Drain decoupled event queues. */
    process_event_queue(&game->events);
    audio_process(&game->audio_queue);

    /* 8. Render into the next buffer and swap. */
    scene_draw(&game->scene_buffers);
}
\end{lstlisting}

\medskip\hrule\medskip

\section{Implementation Checklist}

\begin{enumerate}
\item Prefer \textbf{explicit state} over scattered Boolean flags.
\item Prefer \textbf{command records} when calls must be queued, logged, replayed, undone, or sent over a network.
\item Prefer \textbf{flyweights} for immutable shared data with many instances.
\item Use \textbf{observers} only where dynamic decoupling is worth the loss of static call visibility.
\item Use \textbf{event queues} when sender and receiver should be decoupled in time.
\item Use \textbf{components} to avoid deep inheritance trees and to isolate subsystem behavior.
\item Use \textbf{dirty flags} for derived data that changes less often than it is queried.
\item Use \textbf{object pools} for high-churn fixed-size object families.
\item Use \textbf{packed arrays and active-prefix partitions} for hot per-frame loops.
\item Use \textbf{spatial partitioning} to avoid all-pairs tests.
\end{enumerate}


\clearpage
\section*{Reference Note}
\addcontentsline{toc}{section}{Reference Note}
This handout is derived from a systematic reading of the uploaded source text and its implementation examples. The algorithms have been rewritten in independent C-like pseudocode for engineering study and are intended for lawful technical analysis, architecture review, and implementation planning.

\end{document}
